\begin{figure}[ht]
\centering
\begin{tikzpicture}[x=1cm, y=1cm, font=\small]

% ---- Axes (log-frequency horizontal, weight vertical) ----
\draw[->, thick] (0, 0) -- (10.6, 0) node[right, font=\footnotesize]
  {frequency $f$ (log scale)};
\draw[->, thick] (0, 0) -- (0, 4.4) node[above, font=\footnotesize]
  {blending weight};

% decade ticks: 0.01, 0.1, 1, 10, 100 Hz mapped across 0..10 (2.5 cm/decade)
\foreach \x/\lab in {0/0.01, 2.5/0.1, 5/1, 7.5/10, 10/100} {
  \draw (\x, -0.1) -- (\x, 0.1);
  \node[font=\scriptsize, anchor=north] at (\x, -0.1) {\lab};
}
% weight ticks 0 and 1 (mapped 0 -> 0, 1 -> 3.6)
\foreach \y/\lab in {0/0, 1.8/0.5, 3.6/1} {
  \draw (-0.1, \y) -- (0.1, \y);
  \node[font=\scriptsize, anchor=east] at (-0.1, \y) {\lab};
}

% crossover corner at 1.6 Hz. log10(1.6) = 0.204 -> x = 5 + 2.5*0.204 = 5.51
\draw[thin, enggray, dashed] (5.51, 0) -- (5.51, 3.9);
\node[font=\scriptsize, color=enggray, anchor=south] at (5.51, 3.9)
  {corner $\approx 1.6\,$Hz};

% Accelerometer low-pass weight: 1 at low f, rolling to 0 above corner.
% w_acc(f) = 1 / (1 + (f/fc)^2)^{1/2} -> use smooth S descending
\draw[very thick, engblue]
  (0, 3.55)
  .. controls (2.5, 3.5) and (4.2, 3.35) ..
  (5.51, 1.8)
  .. controls (6.8, 0.55) and (8.0, 0.28) ..
  (10, 0.12);
\node[font=\footnotesize, color=engblue, anchor=west] at (0.2, 3.15)
  {accelerometer (low-pass, gravity reference)};

% Gyroscope high-pass weight: 0 at low f, rising to 1 above corner (mirror image)
\draw[very thick, engred]
  (0, 0.15)
  .. controls (2.5, 0.28) and (4.2, 0.55) ..
  (5.51, 1.8)
  .. controls (6.8, 3.35) and (8.0, 3.5) ..
  (10, 3.58);
\node[font=\footnotesize, color=engred, anchor=east] at (9.8, 3.15)
  {gyroscope (high-pass, integrated rate)};

% crossover point marker
\fill[enggray] (5.51, 1.8) circle (1.8pt);
\node[font=\scriptsize, color=enggray, anchor=south west] at (5.65, 1.85)
  {equal weight $0.5$};

\end{tikzpicture}
\caption[Complementary-filter crossover]{The blending weights of the
complementary filter
$\hat{\theta}_{k+1} = \alpha(\hat{\theta}_{k} + \omega_{k}\Delta t)
+ (1-\alpha)\theta_{\mathrm{acc},k+1}$, plotted against frequency.
Below the crossover corner the accelerometer's gravity-vector
reading (blue) carries the estimate, because it does not drift; above
the corner the integrated gyroscope rate (red) carries it, because
its short-term noise floor is below the accelerometer's
linear-acceleration corruption. For $\alpha = 0.98$ at
$\Delta t = 5\,$ms the corner sits near $1.6\,$Hz, where each sensor
contributes half the estimate. The filter is a one-line fusion rule
that assigns each sensor the frequency band in which it lies least.}
\label{fig:vol01:ch05:complementary-crossover}
\end{figure}
